\documentclass{exam}
\usepackage{../../shah292}

\hypersetup{colorlinks=true, linktoc=section, linkcolor=blue}

\pagestyle{headandfoot}
\firstpageheadrule
\runningheadrule
\firstpageheader{Prof. Shen \\ Differential Equations}{Homework \#3}{Jeevan Shah}
\runningheader{Differential Equations \\ Homework \#3}{}{Shah}
\firstpagefooter{}{\thepage}{}
\runningfooter{ }{\thepage}{ }

\printanswers


\begin{document}
\colorbox{red}{\underline{3.5.9}}
\begin{solution}
    We start by finding the eigenvalues for $A$:
    \begin{align*}
        \det(A-\lambda I) = \begin{vmatrix}
            5 - \lambda & -3 & -2 \\
            8 & -5-\lambda & -4 \\
            -4 & 3 & 3 - \lambda
        \end{vmatrix}
        &\xrightarrow{R_3 \to R_1 + R_3}
        \begin{vmatrix}
            5 - \lambda & -3 & -2 \\
            8 & -5-\lambda & -4 \\
            1 - \lambda & 0 & 1 - \lambda
        \end{vmatrix} \\
        &\xrightarrow{R_1 \to 2R_1 - R_2}
        \begin{vmatrix}
            2 - 2\lambda & -3 & 0 \\
            8 & -5-\lambda & -4 \\
            1 - \lambda & 0 & 1 - \lambda
        \end{vmatrix} \\
        &\xrightarrow{C_1 \to C_1 + 2C_3}
        \begin{vmatrix}
            2 - 2\lambda & -3 & 0 \\
            0 & -5-\lambda & -4 \\
            3 - 3\lambda & 0 & 1 - \lambda
        \end{vmatrix} \\
        &\Rightarrow (\lambda -1)^3 = 0 \Rightarrow \lambda_{1, 2, 3} = 1.
    \end{align*}
    The associated eigenvector is 
    \[
        A - I = \begin{pmatrix}
            4 & -3 & -2 \\
            8 & -6 & -4 \\
            -4 & 3 & 2
        \end{pmatrix} 
        \xrightarrow{}
        \begin{pmatrix}
            4 & -3 & -2 \\
            0 & 0 & 0 \\
            0 & 0 & 0 
        \end{pmatrix} 
        \Rightarrow \vec{u} = \begin{pmatrix}
            1 \\ 0 \\ -2
        \end{pmatrix}.
    \]
    Now since 
    \[
        (A-I)^2 = \mathbf{0}
    \]
    every nonzero vector in $\RR^3$ is a generalized eigenvector of eigenvalue 1. Then,
    \begin{align*}
        e^{tA} = e^t\left(I + t \begin{pmatrix}
            4 & -3 & -2 \\
            8 & -6 & -4 \\
            -4 & 3 & 2
        \end{pmatrix}\right) = \boxed{
            e^t\begin{pmatrix}
                1 + 4t & -3t & -2t \\
                8t & 1-6t & -4t \\
                -4t & 3t & 1 + 2t
            \end{pmatrix}
        }
    \end{align*}
\end{solution}

\colorbox{red}{\underline{3.5.11:}}
\begin{solution}
    We start by finding the eigenvalues of $A$: 
    \begin{align*}
        \det(A - \lambda I) = \begin{vmatrix}
            4 - \lambda & -1 & 2 \\
            1 & 2 - \lambda & 2 \\
            -2 & 2 & 1 - \lambda
        \end{vmatrix}
        &\xrightarrow{R_2 \to R_2 - R_1}
        \begin{vmatrix}
            4 - \lambda & -1 & 2 \\
            \lambda - 3 & 3 - \lambda & 0 \\
            -2 & 2 & 1 - \lambda
        \end{vmatrix} \\
        &\xrightarrow{C_2 \to C_2 + C_1}
        \begin{vmatrix}
            4 - \lambda & 3-\lambda & 2 \\
            \lambda - 3 & 0 & 0 \\
            -2 & 0 & 1 - \lambda
        \end{vmatrix} \\
        &\Rightarrow (\lambda - 3)^2(1 - \lambda) \Rightarrow \lambda_{1, 2} = 3 \text{ and } \lambda_3 = 1.
    \end{align*}    
    The corresponding eigenvectors are 
    \begin{align*}
        A - 3I &= \begin{pmatrix}
            1 & -1 & 2 \\
            1 & -1 & 2 \\
            -2 & 2 & -2 
        \end{pmatrix}
        \to
        \begin{pmatrix}
            1 & -1 & 2 \\
            0 & 0 & 0 \\
            0 & 0 & 0
        \end{pmatrix}
        \Rightarrow \vec{u_1} = \begin{pmatrix}
            1 \\ 1 \\ 0
        \end{pmatrix} \\
        A - I &= \begin{pmatrix}
            3 & -1 & 2 \\
            1 & 1 & 2 \\
            -2 & 2 & 0 
        \end{pmatrix}
        \to \begin{pmatrix}
            1 & 1 & 2 \\
            3 & -1 & 2 \\
            -2 & 2 & 0
        \end{pmatrix}
        \to \begin{pmatrix}
            1 & 1 & 2 \\
            0 & -4 & -4 \\
            0 & 4 & 4
        \end{pmatrix}
        \to \begin{pmatrix}
            1 & 1 & 2 \\
            0 & -4 & -4 \\
            0 & 0 & 0
        \end{pmatrix}
        \Rightarrow \vec{u_3} = \begin{pmatrix}
            -1 \\ -1 \\ 1
        \end{pmatrix}.
    \end{align*}
    We must also find a generalized eigenvector since $\lambda = 3$ has a multiplicity of $2$ but the dimension of its eigenspace is $1 < 2$. We can find the generalized eigenvector by solving the equation
    \[
        0 = (A-3I)^2\vec{w} = \begin{pmatrix}
            -4 & 4 & -4 \\
            -4 & 4 & -4 \\
            4 & -4 & 4
        \end{pmatrix}\vec{w}
        \to \begin{pmatrix}
            -1 & 1 & 1 \\
            0 & 0 & 0 \\
            0 & 0 & 0
        \end{pmatrix}
        \Rightarrow \vec{w} = \begin{pmatrix}
            -1 \\ 0 \\ 1
        \end{pmatrix}.
    \]
    Then, 
    \[
        P = \begin{pmatrix}
            \vec{u_1} & \vec{u_3} & \vec{w}
        \end{pmatrix} = \begin{pmatrix}
            -1 & 1 & -1 \\
            -1 & 1 & 0 \\
            1 & 0 & 1
        \end{pmatrix}
    \]
    so the Jordan form is 
    \[
        J = \begin{pmatrix}
            1 & 0 & 0 \\
            0 & 3 & 1 \\
            0 & 0 & 3
        \end{pmatrix}.
    \]
    For the $2 \times 2$ Jordan block $J_2(3) = 3I + N$ where 
    \[
        N = \begin{pmatrix}
            0 & 1 \\
            0 & 0
        \end{pmatrix},
        \quad N^2 = 0.
    \]
    Also 
    \[
        e^{tJ_{2}(3)} = e^{t(3I+N)} = e^{3t}e^{tN} = e^{3t}(I + tN) = e^{3t}\begin{pmatrix}
            1 & t \\ 0 & 1
        \end{pmatrix},
    \]
    so 
    \[
        e^{tJ} = \begin{pmatrix}
            e^t & 0 & 0 \\
            0 & e^{3t} & te^{3t} \\
            0 & 0 & e^{3t}
        \end{pmatrix}.
    \]
    Since $e^{tA} = Pe^{tJ}P^{-1}$, we have 
    \[
        \boxed{
            e^{tA} = e^{3t}\begin{pmatrix}
                2 - t & t - 1 & 1 \\
                1 - t & t & 1 \\
                -1 & 1 & 0
            \end{pmatrix}
            + e^t \begin{pmatrix}
                -1 & 1 & -1 \\
                -1 & 1 & -1 \\
                0 & -1 & 1
            \end{pmatrix}
        }
    \]
\end{solution}

\colorbox{red}{\underline{4.5.2:}}
\begin{solution}
    We start by finding $e^{tA}$. The eigenvalues of $A$ are given by the characteristic polynomial so 
    \[
        \lambda^2 - 10\lambda + 16 \Rightarrow (\lambda -8)(\lambda - 2) = 0 \Rightarrow \lambda_1 = 8, \lambda_2 = 2.
    \]
    The associated eigenvectors are 
    \begin{align*}
        A - 8I &= \begin{pmatrix}
            -3 & 3 \\
            3 & -3
        \end{pmatrix}
        \to \begin{pmatrix}
            -1 & 1 \\
            0 & 0
        \end{pmatrix}
        \Rightarrow \vec{u_1} = \begin{pmatrix}
            1 \\ 0
        \end{pmatrix} \\
        A - 2I &= \begin{pmatrix}
            3 & 3 \\
            3 & 3
        \end{pmatrix}
        \to \begin{pmatrix}
            1 & 1 \\ 
            0 & 0 
        \end{pmatrix}
        \Rightarrow \vec{u_2} = \begin{pmatrix}
            -1 \\ 1
        \end{pmatrix}.
    \end{align*}
    Then, 
    \begin{align*}
        V = \begin{pmatrix}
            1 & -1 \\
            1 & 1
        \end{pmatrix} 
        \Rightarrow V^{-1} = \frac{1}{2}\begin{pmatrix}
            1 & 1 \\
            -1 & 1
        \end{pmatrix}
        \Rightarrow e^{tA} &= \begin{pmatrix}
            1 & -1 \\
            1 & 1 
        \end{pmatrix}
        \begin{pmatrix}
            e^{8t} & 0 \\
            0 & e^{2t}
        \end{pmatrix}
        \frac{1}{2}
        \begin{pmatrix}
            1 & 1 \\
            -1 & 1
        \end{pmatrix}\\
        &= \frac{1}{2}\begin{pmatrix}
            e^{8t} + e^{2t} & e^{8t} - e^{2t} \\
            e^{8t} - e^{2t} & e^{8t} + e^{2t}
        \end{pmatrix}\vec{x}.
    \end{align*}
    So, 
    \begin{align*}
        \vec{x}'(t) = A\vec{x}(t) + \vec{g}(t) \Rightarrow (e^{-tA}\vec{x})' = e^{-tA}\vec{g} &\Rightarrow e^{-tA}\vec{x} - e^{-t_0 A}\vec{x}_0 = \int_{t_0}^{t}e^{sA}\vec{g}(s)\,\mathrm{d}s \\
        &\Rightarrow \vec{x}(t) = e^{t - t_0}\vec{x_0} + \int_{t_0}^{t}e^{(t-s)A}\vec{g}(s)\,\mathrm{d}s.
    \end{align*}
    But, 
    \[
        \vec{g}(s) = e^{-8s}(s, s) = se^{-8s}(1, 1) 
    \]
    and 
    \[
        e^{(t-s)A}(1, 1) = e^{8(t-s)}(1, 1)
    \]
    since $8$ is an eigenvalue of $A$. So 
    \begin{align*}
        \int_{0}^{t} e^{(t-s)A}e^{-8s}(s, s)^{T}\,\mathrm{d}s &= \int_{0}^{t}se^{-8s}e^{8(t-s)}\,\mathrm{d}s \,(1, 1)^{T} \\
        &= e^{8t}\left(\int_{0}^{t}se^{-16s}\,\mathrm{d}s\right)\begin{pmatrix}
            1 \\ 1
        \end{pmatrix} \\
        &= e^{8t}\left(\frac{1}{256} - e^{-16t}\left(\frac{t}{16} + \frac{1}{256}\right)\right)\begin{pmatrix}
            1 \\ 1
        \end{pmatrix}\\ 
        &= \frac{e^{8t} - (16t + 1)e^{-8t}}{256}\begin{pmatrix}
            1 \\ 1
        \end{pmatrix}
    \end{align*}
    Since 
    \[
        e^{tA}{\vec{x}_0}  = \begin{pmatrix}
            \frac{e^{8t} - e^{2t}}{2} \\
            \frac{e^{8t} + e^{2t}}{2}
        \end{pmatrix}
    \]
    we have 
    \[
        \boxed{
            \vec{x}(t) = \begin{pmatrix}
            \frac{e^{8t} - e^{2t}}{2} \\
            \frac{e^{8t} + e^{2t}}{2}
        \end{pmatrix} + \frac{e^{8t} - (16t + 1)e^{-8t}}{256}\begin{pmatrix}
            1 \\ 1
        \end{pmatrix}
        }
    \]
\end{solution}

\colorbox{red}{\underline{4.5.8:}}
\begin{solution}
    \begin{parts}
    \part We must solve the system $\vec{v}(x, y) = \vec{0}$ to the equilibrium points. This is equivalent to 
    \begin{align*}
        0 &= (x + y)(x - y - 1) \Rightarrow x + y = 0 \quad\text{or}\quad x - y - 1 = 0 \\
        0 &= (x + y - 2)(x - y + 1) \Rightarrow x + y - 2 = 0 \quad\text{or}\quad x - y + 1 = 0.
    \end{align*}
    The equilibrium points will lie at the intersection between pairs of these equations. One can see that the only consistent pairs are 
    \begin{align*}
        &x + y = 0 \quad\text{and}\quad x - y + 1 = 0 \Rightarrow (x, y) = \left(-\frac{1}{2}, \frac{1}{2}\right) \\
        &x - y - 1 = 0 \quad\text{and}\quad x + y - 2 = 0 \Rightarrow (x, y) = \left(\frac{3}{2}, \frac{1}{2}\right).
    \end{align*}
    Linearizing we see that 
    \[
        f(x, y) = (x + y)(x - y - 1) \quad\text{and}\quad g(x, y) = (x + y - 2)(x - y +1)
    \]
    so 
    \[
        D_v(x, y) = \begin{pmatrix}
            f_x & f_y \\
            g_x & g_y
        \end{pmatrix} 
        = \begin{pmatrix}
            2x - 1 & -2y - 1 \\
            2x - 1 & -2y + 3
        \end{pmatrix}.
    \]
    Then 
    \begin{align*}
        D_v\left(-\frac{1}{2}, \frac{1}{2}\right) &= \begin{pmatrix}
            -2 & -2 \\
            -2 & 2
        \end{pmatrix}
        \Rightarrow \det(D_v) = - 8 < 0 \\
        D_v\left(\frac{3}{2}, \frac{1}{2}\right) &= \begin{pmatrix}
            2 & -2 \\
            2 & 2
        \end{pmatrix}
        \Rightarrow \lambda_{1, 2} = 2 \pm 2i.
    \end{align*}
    It follows that \textbf{both equilibrium points are unstable.}
    
    \part We know from part $(a)$ that the equilibrium points lie at 
    \[
        \left(-\frac{1}{2}, \frac{1}{2}\right) \quad\text{and}\quad \left(\frac{3}{2}, \frac{1}{2}\right).
    \]
    Then, 
    \[
        f(x, y) = (x + y - 2)(x - y + 1) \quad\text{and}\quad g(x, y) = (x + y)(x - y - 1)
    \]
    so 
    \[
        D_v(x, y) = \begin{pmatrix}
            f_x & f_y \\
            g_x & g_y
        \end{pmatrix} 
        = \begin{pmatrix}
            2x - 1 & -2y + 3 \\
            2x - 1 & -2y - 1
        \end{pmatrix}.
    \]
    Then 
    \begin{align*}
        D_v\left(-\frac{1}{2}, \frac{1}{2}\right) &= \begin{pmatrix}
            -2 & 2 \\
            -2 & -2
        \end{pmatrix}
        \Rightarrow \lambda_{1, 2} = -2 \pm 2i \\
        D_v\left(\frac{3}{2}, \frac{1}{2}\right) &= \begin{pmatrix}
            2 & 2 \\
            2 & -2
        \end{pmatrix}
        \Rightarrow \det(D_v) = -8 < 0
    \end{align*}
    It follows that $\left(-\frac{1}{2}, \frac{1}{2}\right)$ \textbf{is stable} and $\left(\frac{3}{2}, \frac{1}{2}\right)$ \textbf{is unstable}.
    
\end{parts}
\end{solution}

\colorbox{red}{\underline{4.5.9:}}
\begin{solution}
    \begin{parts}
        \part Just like in $(4.5.8)$ we start by solving $\vec{v}(x, y) = \vec{0}$. This is equivalent to 
        \begin{align*}
            x - xy &= x(1-y) = 0 \Rightarrow x = 0 \quad\text{or}\quad y = 1  \\
            y + 2xy &= y(1 + 2x) \Rightarrow x = -\frac{1}{2} \quad\text{or}\quad y = 0.
        \end{align*}
        So the equilibrium points are $(0, 0)$ and $\left(-\frac{1}{2}, 1\right)$. Then, 
        \[
            f(x, y) = x - xy \quad\text{and}\quad g(x, y) = y + 2xy 
        \]
        so
        \[
            D_v(x, y) = \begin{pmatrix}
                1 - y & -x \\
                2y & 1 + 2x
            \end{pmatrix}.
        \]
        Thus, 
        \begin{align*}
            D_v(0, 0) &= \begin{pmatrix}
                1 & 0 \\
                0 & 1
            \end{pmatrix}
            \Rightarrow \lambda_{1, 2} = 1 \\
            D_v\left(-\frac{1}{2}, 1\right) &= \begin{pmatrix}
                0 & \frac{1}{2} \\
                2 & 0 
            \end{pmatrix}
            \Rightarrow \det(D_v) = - 1 < 0.
        \end{align*}
        It follows that \textbf{both equilibrium points are unstable}. 

        \part From $(a)$ we know that $(0, 0)$ and $\left(-\frac{1}{2}, 1\right)$ are equilibrium points. Then 
        \[
            f(x, y) = y + 2xy \quad\text{and}\quad g(x, y) = x - xy 
        \]
        so
        \[
            D_v(x, y) = \begin{pmatrix}
                2y & 1 + 2x \\
                1 - y & -x 
            \end{pmatrix}.
        \]
        Thus, 
        \begin{align*}
            D_v(0, 0) &= \begin{pmatrix}
                0 & 1 \\
                1 & 0
            \end{pmatrix}
            \Rightarrow \det(D_v) = -1 < 0 \\
            D_v\left(-\frac{1}{2}, 1\right) &= \begin{pmatrix}
                2 & 0 \\
                0 & \frac{1}{2}
            \end{pmatrix}
            \Rightarrow \lambda_{1} = 2, \lambda_{2} = \frac{1}{2} \Rightarrow \mathfrak{R}(\lambda_1) > 0, \mathfrak{R}(\lambda_2) > 0.
        \end{align*}
        It follows that \textbf{both equilibrium points are unstable}. 
    \end{parts}
\end{solution}
\end{document}